% !TeX root = main.tex
% Synchronized with 00_摘要.md v0.5; Tsai format.
\begin{abstracten}
\begin{spacing}{1.5}

Mobile manipulators must coordinate navigation and manipulation under changing obstacles, reachability constraints, and arm tracking lag. This thesis develops a control pipeline combining predictive navigation, whole-body coordination, and visual guidance. Navigation uses geometric model predictive control with \(SE(2)\) errors, predictive barrier constraints, and a raw-laser-scan protection layer. Coordinated control incorporates arm position setpoints and tracking dynamics to jointly optimize base and arm commands.

In simulations, the navigation method reached the goal on all 83 routes jointly evaluated across three methods, with contact on one route. Three paired drawer experiments showed task-time reductions of 13.38--14.05\% relative to a strict stationary-base strategy, with both strategies meeting the original criteria. A separate three-pair study showed lower end-effector position errors over the full control window and fewer arm-command reversals with five-step prediction than with one-step prediction. Visual bottle pick-and-place and door operation achieved specified case-level results. The door trial used ideal grasp binding and non-real-time computation. Complete residential retrieval, dynamic whole-body avoidance, and real-time hardware execution remain to be validated.

\keyworden{mobile manipulators, geometric model predictive control, control barrier functions, whole-body coordination, manipulation on the move, visual guidance}
\end{spacing}
\end{abstracten}
